\documentclass[10pt,a4paper]{amsart}
\usepackage{amsmath,amssymb,amsthm,graphicx}
\usepackage{verbatim}
\newenvironment{reference}{\comment}{\endcomment}
%\newenvironment{reference}{}{}
\newenvironment{slogan}{\comment}{\endcomment}
\newenvironment{history}{\comment}{\endcomment}
% For commutative diagrams we use Xy-pic
\usepackage[all]{xy}

% We use 2cell for 2-commutative diagrams.
\xyoption{2cell}
% Theorem environments.
%
\theoremstyle{plain}
\newtheorem{theorem}[subsection]{Theorem}
\newtheorem{proposition}[subsection]{Proposition}
\newtheorem{lemma}[subsection]{Lemma}

\theoremstyle{definition}
\newtheorem{definition}[subsection]{Definition}
\newtheorem{example}[subsection]{Example}
\newtheorem{exercise}[subsection]{Exercise}
\newtheorem{situation}[subsection]{Situation}

\theoremstyle{remark}
\newtheorem{remark}[subsection]{Remark}
\newtheorem{remarks}[subsection]{Remarks}

\numberwithin{equation}{subsection}

% Macros
%
\def\lim{\mathop{\mathrm{lim}}\nolimits}
\def\colim{\mathop{\mathrm{colim}}\nolimits}
\def\Spec{\mathop{\mathrm{Spec}}}
\def\Hom{\mathop{\mathrm{Hom}}\nolimits}
\def\Ext{\mathop{\mathrm{Ext}}\nolimits}
\def\SheafHom{\mathop{\mathcal{H}\!\mathit{om}}\nolimits}
\def\SheafExt{\mathop{\mathcal{E}\!\mathit{xt}}\nolimits}
\def\Sch{\mathit{Sch}}
\def\Mor{\mathop{\mathrm{Mor}}\nolimits}
\def\Ob{\mathop{\mathrm{Ob}}\nolimits}
\def\Sh{\mathop{\mathit{Sh}}\nolimits}
\def\NL{\mathop{N\!L}\nolimits}
\def\CH{\mathop{\mathrm{CH}}\nolimits}
\def\proetale{{pro\text{-}\acute{e}tale}}
\def\etale{{\acute{e}tale}}
\def\QCoh{\mathit{QCoh}}
\def\Ker{\mathop{\mathrm{Ker}}}
\def\Im{\mathop{\mathrm{Im}}}
\def\Coker{\mathop{\mathrm{Coker}}}
\def\Coim{\mathop{\mathrm{Coim}}}

% Boxtimes
%
\DeclareMathSymbol{\boxtimes}{\mathbin}{AMSa}{"02}

%
% Macros for moduli stacks/spaces
%
\def\QCohstack{\mathcal{QC}\!\mathit{oh}}
\def\Cohstack{\mathcal{C}\!\mathit{oh}}
\def\Spacesstack{\mathcal{S}\!\mathit{paces}}
\def\Quotfunctor{\mathrm{Quot}}
\def\Hilbfunctor{\mathrm{Hilb}}
\def\Curvesstack{\mathcal{C}\!\mathit{urves}}
\def\Polarizedstack{\mathcal{P}\!\mathit{olarized}}
\def\Complexesstack{\mathcal{C}\!\mathit{omplexes}}
% \Pic is the operator that assigns to X its picard group, usage \Pic(X)
% \Picardstack_{X/B} denotes the Picard stack of X over B
% \Picardfunctor_{X/B} denotes the Picard functor of X over B
\def\Pic{\mathop{\mathrm{Pic}}\nolimits}
\def\Picardstack{\mathcal{P}\!\mathit{ic}}
\def\Picardfunctor{\mathrm{Pic}}
\def\Deformationcategory{\mathcal{D}\!\mathit{ef}}
\usepackage{cleveref}
\crefname{originallemma}{original lemma}{original lemmas}
\Crefname{originallemma}{Original lemma}{Original lemmas}
\makeatletter
\newlabel{lemma-top-chern-class}{{0FA9}{source}{}{sourcepointer-1}{}}
\newlabel{lemma-top-chern-class@cref}{{[originallemma][0][]0FA9}{[1][0][]source}{}{}{}}
\makeatother
\crefname{originalremark}{original remark}{original remarks}
\Crefname{originalremark}{Original remark}{Original remarks}
\makeatletter
\newlabel{remark-top-chern-class}{{0FA4}{source}{}{sourcepointer-2}{}}
\newlabel{remark-top-chern-class@cref}{{[originalremark][0][]0FA4}{[1][0][]source}{}{}{}}
\makeatother
\crefname{originalreference}{original reference}{original references}
\Crefname{originalreference}{Original reference}{Original references}
\makeatletter
\newlabel{situation-setup}{{02QL}{source}{}{sourcepointer-3}{}}
\newlabel{situation-setup@cref}{{[originalreference][0][]02QL}{[1][0][]source}{}{}{}}
\makeatother
\title{The class of the total space of a subbundle as the cap of a top Chern class}
\author{The Stacks Project Authors}
\date{stacks-project:chow.tex:local-snapshot-20261012; GNU FDL 1.2 or later}
\begin{document}
\maketitle
\noindent\textit{Editorial note.} All delivered mathematical text is the author's own, copied from the pinned Stacks Project chapter file chow.tex; no statement or proof step was rewritten, reordered or removed.\par\medskip
\noindent\textit{Editorial note.} Preamble provenance: the pinned chapter file opens with \textbackslash{}input\{preamble\} and carries no preamble text of its own, so the project's single Stacks-wide preamble staged at sources/stacks-chow/preamble.tex is excerpted verbatim as three spans, lines 8-12 (the verbatim package together with the project's own reference, slogan and history display-environment definitions, including the author's commented-out older form of the reference environment), lines 14-18 (Xy-pic setup) and lines 154-217 (the native theorem declarations and mathematical macros). The lines 8-12 span is required because a counted statement or its surrounding text may use those display environments, which the project typesets as comments. Nothing was added, removed or redefined; the wrapper class and the amsmath/amssymb/amsthm/graphicx base packages are supplied by the assembler.\par\medskip
\noindent\textit{Editorial note.} Same-chapter pointers are located at their \textbackslash{}label lines inside the pinned chapter file; cross-chapter pointers are located in the project's own tag list staged in the snapshot (sources/stacks-chow/tags.txt), at the verbatim lines carrying their canonical tags.\par\medskip
\noindent\textit{Allowed standard prerequisites (editorial).}\begin{itemize}
\item The bivariant top Chern class attached to a finite locally free module. Same chapter chow.tex, label remark-top-chern-class, invoked at proof line 8861. Fixes what the class c\_top of the quotient module is and justifies verifying the formula first for modules of constant rank.
\item The geometric case of the top Chern class for a regular section. Same chapter chow.tex, label lemma-top-chern-class, invoked at proof line 8892. Supplies the identification of the class of the vanishing scheme of a section cut out by a regular sequence with the cap of the top Chern class, which is exactly the formula for a single integral class.
\end{itemize}\medskip
\ifdefined\hypertarget\hypertarget{sourcepointer-1}{}\fi
\noindent\textit{Source reference.} Original lemma 0FA9: sources/stacks-chow/pinned-original-source.tex, source lines 8750--8750. Navigation pointer for \textbackslash{}ref\{lemma-top-chern-class\} used inside the counted span (line 8892); the argument that consumes it is retained in the span.\par\medskip
\ifdefined\hypertarget\hypertarget{sourcepointer-2}{}\fi
\noindent\textit{Source reference.} Original remark 0FA4: sources/stacks-chow/pinned-original-source.tex, source lines 7677--7677. Navigation pointer for \textbackslash{}ref\{remark-top-chern-class\} used inside the counted span (line 8862); the argument that consumes it is retained in the span.\par\medskip
\ifdefined\hypertarget\hypertarget{sourcepointer-3}{}\fi
\noindent\textit{Source reference.} Original reference 02QL: sources/stacks-chow/pinned-original-source.tex, source lines 1372--1372. Navigation pointer for \textbackslash{}ref\{situation-setup\} used inside the counted span (line 8840); the argument that consumes it is retained in the span.\par\medskip
\begin{lemma}
\label{lemma-easy-virtual-class}
Let $(S, \delta)$ be as in Situation \ref{situation-setup}. Let $X$
be a scheme locally of finite type over $S$. Let
$$
0 \to \mathcal{N}' \to \mathcal{N} \to \mathcal{E} \to 0
$$
be a short exact sequence of finite locally free $\mathcal{O}_X$-modules.
Consider the closed embedding
$$
i :
N' = \underline{\Spec}_X(\text{Sym}((\mathcal{N}')^\vee))
\longrightarrow
N = \underline{\Spec}_X(\text{Sym}(\mathcal{N}^\vee))
$$
For $\alpha \in \CH_k(X)$ we have
$$
i_*(p')^*\alpha = p^*(c_{top}(\mathcal{E}) \cap \alpha)
$$
where $p' : N' \to X$ and $p : N \to X$ are the structure morphisms.
\end{lemma}

\begin{proof}
Here $c_{top}(\mathcal{E})$ is the bivariant class defined in
Remark \ref{remark-top-chern-class}. By its very definition, in
order to verify the formula, we may assume that $\mathcal{E}$
has constant rank. We may similarly assume $\mathcal{N}'$ and
$\mathcal{N}$ have constant ranks, say $r'$ and $r$, so
$\mathcal{E}$ has rank $r - r'$ and
$c_{top}(\mathcal{E}) = c_{r - r'}(\mathcal{E})$.
Observe that $p^*\mathcal{E}$ has a canonical section
$$
s \in \Gamma(N, p^*\mathcal{E}) = \Gamma(X, p_*p^*\mathcal{E}) =
\Gamma(X, \mathcal{E} \otimes_{\mathcal{O}_X} \text{Sym}(\mathcal{N}^\vee)
\supset \Gamma(X, \SheafHom(\mathcal{N}, \mathcal{E}))
$$
corresponding to the surjection $\mathcal{N} \to \mathcal{E}$ given
in the statement of the lemma. The vanishing scheme of this section
is exactly $N' \subset N$. Let $Y \subset X$ be an integral closed
subscheme of $\delta$-dimension $n$. Then we have
\begin{enumerate}
\item $p^*[Y] = [p^{-1}(Y)]$ since $p^{-1}(Y)$ is integral of
$\delta$-dimension $n + r$,
\item $(p')^*[Y] = [(p')^{-1}(Y)]$ since $(p')^{-1}(Y)$ is integral of
$\delta$-dimension $n + r'$,
\item the restriction of $s$ to $p^{-1}Y$ has vanishing scheme
$(p')^{-1}Y$ and the closed immersion $(p')^{-1}Y \to p^{-1}Y$
is a regular immersion (locally cut out by a regular sequence).
\end{enumerate}
We conclude that
$$
(p')^*[Y] = c_{r - r'}(p^*\mathcal{E}) \cap p^*[Y]
\quad\text{in}\quad \CH_*(N)
$$
by Lemma \ref{lemma-top-chern-class}. This proves the lemma.
\end{proof}

\end{document}
